\chapter{Bevezetés}

A digitalizáció elterjedésével egyre nagyobb szerepet kap az adatok elektronikus formában történő gyűjtése, feldolgozása és továbbítása. Számos területen szükség van olyan megoldásokra, amelyek lehetővé teszik az adatok egységes szerkezetben történő rögzítését, valamint biztosítják, hogy a különböző rendszerek között továbbított információk megfelelő formátumúak legyenek. Az XML (Extensible Markup Language) erre egy széles körben használt lehetőség, aminek a segítségével az adatok egy hierarchikus szerkezetben, egyszerűen értelmezhető módon, és géppel is feldolgozható formában tárolhatóak.

Az XML dokumentumok szerkezetének és tartalmának meghatározására többek között az XSD (XML Schema Definition) szolgál. Egy XSD séma megadja, hogy milyen elemekből épülhet fel egy XML dokumentum, és az egyes elemek típusára, kötelezőségére, előfordulási számára és megengedett értékeire vonatkozó megkötéseket is tartalmazhat. Ezzel lehetővé válik, hogy az adatok rögzítése előre meghatározott szabályok szerint történjen. Ugyanakkor egy ilyen séma alapján történő adatbevitel a felhasználók számára nem feltétlenül egyszerű feladat, ha nem rendelkeznek XML-lel kapcsolatos ismeretekkel. Az XML dokumentumok kézi létrehozása és szerkesztése összetettebb adatszerkezetek esetén időigényes, és könnyű hibát véteni a folyamat alatt.

A probléma egy megoldása egy olyan grafikus felület biztosítása, amely az XSD séma alapján automatikusan létrehozza az adatok beviteléhez szükséges űrlapot. Ebben az esetben a felhasználónak nem kell ismernie az XML dokumentumok felépítését vagy az azokhoz kapcsolódó szabályokat, hanem egy hagyományos böngészőben megjeleníthető űrlapon keresztül adhatja meg a szükséges információkat. Ezzel könnyűvé válik az adatgyűjtés, miközben az előre meghatározott séma által az adatok szerkezete is megfelelő marad. Egy ilyen alkalmazás megvalósítása azonban nem csupán az űrlapmezők létrehozását jelenti, figyelembe kell venni az XSD által meghatározott összetett adatszerkezeteket, az egyes adattípusokat, valamint a különböző korlátozásokat is.

A szakdolgozat célja egy olyan webalkalmazás megtervezése és megvalósítása, amely lehetőséget biztosít XSD sémák feltöltésére, szerkesztésére, valamint az azok alapján történő dinamikus űrlapgenerálásra. Az alkalmazás segítségével a séma létrehozója meghatározhatja az adatbevitel szerkezetét, majd az elkészült űrlapot más felhasználók számára is elérhetővé teheti. Az űrlap kitöltése során a felhasználónak elég a megfelelő mezőkben megadnia az adatokat, amelyeket egy használható XML dokumentummá állít össze a program.